\section{An obstruction in Steenrod modules}\label{sec:algebra}

The mapping-cone construction in Section~\ref{sec:cones} supplies
a nonzero tensor on a weight summand. In our application, its detecting
functional also lifts from an annihilated functional at the preceding level.
The obstruction below uses both inputs. Cartier halving of the lifted
functional gives a lower bound on the $d_0$-exponent. An
even-degree Tor injection then lets us halve the tensor without losing
its nonvanishing. The exponent bound controls the change of ideal at
each step.

Throughout this section $k\geq1$,
\[
 P_k=\F_2[x_1,\ldots,x_k],\qquad
 D_k=P_k^{\GL_k(\F_2)}=\F_2[d_0,\ldots,d_{k-1}],\qquad
 |d_j|=2^k-2^j,
\]
and $d_k=1$.  The invariants are normalized by
\begin{equation}\label{eq:algebra-dickson-polynomial}
 P(U)=\prod_{a\in\langle x_1,\ldots,x_k\rangle_{\F_2}}(U+a)
     =\sum_{j=0}^k d_jU^{2^j}.
\end{equation}
For $n\geq0$ we use the module
$M_k(n)=\Sigma^n d_0^nD_k$ of
Lemma~\ref{lem:weight-modules}.  A monomial in $M_k(n)$ means its
underlying polynomial monomial, followed by the grading shift.
In particular,
\begin{equation}\label{eq:algebra-degree-parity}
 \left|\Sigma^n d_0^c\prod_{j=1}^{k-1}d_j^{b_j}\right|
 =n+c(2^k-1)+\sum_{j=1}^{k-1}b_j(2^k-2^j)
 \equiv n+c\pmod2.
\end{equation}
An annihilated functional on a graded $A$-module $M$ is a homogeneous
linear functional $\lambda$ such that $\lambda(IM)=0$.

The Steenrod indecomposables of the ideals \(d_0^pD_k\) have been
studied by Giambalvo and Peterson \cite{GiambalvoPeterson2001}.
Here we first take coinvariants for the exterior algebra of Milnor
primitives, and retain the least allowed \(d_0\)-exponent throughout
the subsequent degree-halving argument. Both the lifted functional
and the tensor nonvanishing in the next statement will be needed.

The tensor product \(I\otimes_A M\) uses the right \(A\)-module
structure on \(I\); its balancing relation is
\((\alpha\beta)\otimes z=\alpha\otimes(\beta z)\), for
\(\alpha\in I\) and \(\beta\in A\).

\begin{lemma}[Algebra obstruction]\label{lem:algebra-obstruction}
Let $k\geq1$, $n=e+1$ with $e\geq0$, and $m\geq2$.  Let $v$ be
a degree-$m$ monomial of $M_k(n)$.  Suppose that its underlying
monomial is detected by an annihilated functional on $M_k(e)$ in
degree $m-1$; equivalently, $v$ is detected by the shifted restriction
of that functional along the ideal inclusion.
If
\begin{equation}\label{eq:obstruction-tensor}
 \Sq^m\otimes v\neq0
       \qquad\text{in } I\otimes_AM_k(n),
\end{equation}
then $m$ is a power of two and $n=k=1$.

Without the functional hypothesis, the tensor in
\eqref{eq:obstruction-tensor} is still zero whenever $n>0$ and
$m\geq3$ is odd.
\end{lemma}

The two halving arguments start at different exponent floors: the
functional is defined on $M_k(n-1)$, whereas the tensor lies over
$M_k(n)$. We first compute the Milnor coinvariants and Cartier map,
then obtain the functional's exponent bound. The Tor injection and
the remaining square calculations will turn that bound into the
asserted tensor obstruction.

\subsection{Milnor primitives and coinvariants}

Let $E\subset A$ be the exterior sub-Hopf algebra on the Milnor
primitives $q_i$, indexed by
$|q_i|=2^{i+1}-1$ for $i\geq0$.  Put $I_E=E^{>0}$.

\begin{lemma}\label{lem:algebra-milnor-quotient}
There is a quotient homomorphism
\[
 V:A\longrightarrow A,\qquad
 V(\Sq^{2a})=\Sq^a,\qquad V(\Sq^{2a+1})=0,
\]
where the target is regarded as having doubled degrees.  Its kernel is
\[
 K=I_EA=AI_E.
\]
Consequently $H_0(E,M)=M/I_EM$ is naturally a module over
$A/K$, for every $A$-module $M$.
\end{lemma}

\begin{proof}
Milnor's polynomial dual and coproduct description
\cite[\S5, Theorem~3, and Appendix~1]{Milnor1958} gives
$A_*=\F_2[\xi_1,\xi_2,\ldots]$. It implies that restriction to
$E_*$ is the quotient by the squares of these generators: their
images are the exterior generators dual to the Milnor primitives.
Squaring is an injective Hopf algebra
homomorphism on $A_*$; its degreewise dual is the indicated map $V$.
The formulas on the squares follow by pairing with the powers of
$\xi_1$.

For completeness, the kernel can be read directly from the Milnor
coproduct.  After restricting its first factor to $E_*$, that coproduct
becomes
\[
 \xi_i\longmapsto \bar\xi_i\otimes1+1\otimes\xi_i.
\]
The other terms have an even positive exponent in the first factor and
therefore vanish.  A polynomial has trivial coaction under this map
exactly when it is a square: the coefficients of the single generators
$\bar\xi_i$ are its ordinary partial derivatives, and their simultaneous
vanishing in characteristic two is equivalent to every exponent being
even.  Conversely, a square has trivial coaction since
$\bar\xi_i^2=0$.  Thus the annihilator of $I_EA$ is exactly the subalgebra
of squares, which also annihilates $\ker V$.  Degreewise duality proves
$I_EA=\ker V$.  Conjugation preserves $E$ and the subalgebra of squares
in $A_*$, so it also gives $AI_E=\ker V$.  Finally
$KM=I_EM$, using $K=I_EA$ and $AM=M$, and hence the stated quotient
module is well defined.
\end{proof}

\begin{lemma}\label{lem:milnor-coinvariants}
The first $k$ Milnor primitives act on $D_k$ by
\begin{equation}\label{eq:milnor-action}
 q_i=d_0\frac{\partial}{\partial d_{i+1}}\quad(0\leq i<k-1),
 \qquad
 q_{k-1}=d_0\sum_{j=0}^{k-1}d_j\frac{\partial}{\partial d_j}.
\end{equation}
Every later primitive is a linear combination of these with coefficients
that are squares in $D_k$.  All ideals $d_0^pD_k$, $p\geq0$, are $A$-submodules.

The quotient $H_0(E,d_0^pD_k)$ has a monomial basis consisting of all
monomials of $d_0$-exponent $p$, together with
\begin{equation}\label{eq:coinvariants}
 d_0^c\prod_{j=1}^{k-1}d_j^{b_j},\qquad
 c>p,\quad c\equiv k\pmod2,\quad b_j\text{ odd for every }j\geq1.
\end{equation}
\end{lemma}

\begin{proof}
The $q_i$ are derivations and $q_i x_s=x_s^{2^{i+1}}$.
The initial formulas on Dickson generators appear in
\cite[Corollary~2.3(b)]{Wilkerson1983} and are also recorded, with
indexing shifted by one, in \cite[Proposition~2.3]{Sum2020}; we give the
calculation.  Apply $q_i$ to $P(x_s)=0$.  Only the linear power of $x_s$
contributes a derivative, so
\[
 \sum_{j=0}^{k-1}q_i(d_j)x_s^{2^j}=d_0x_s^{2^{i+1}}.
\]
This identity holds on the whole binary span of the $x_s$, by additivity.
A linearized polynomial of degree at most $2^{k-1}$ is uniquely
determined by its values on that span: the difference of two such
polynomials cannot have $2^k$ distinct roots unless it is zero.
For $i<k-1$ this gives the partial-derivative formula.  For $i=k-1$,
substituting $x_s^{2^k}=\sum_{j<k}d_jx_s^{2^j}$ gives the Euler formula.

For $t\geq0$, raising the same root equation to the power $2^{t+1}$
gives, as an equality of derivations on $P_k$,
\[
 q_{k+t}=\sum_{j=0}^{k-1}d_j^{2^{t+1}}q_{j+t}.
\]
Induction on $t$ expresses every later primitive in the first $k$ with
square coefficients.  These coefficients belong to $D_k$.

Since $d_0$ is the product of the nonzero binary linear forms and the
total square of such a form $a$ is $a+a^2$, the total square of $d_0$
is divisible by $d_0$.  Its quotient is invariant and therefore belongs
to $D_k$.  The Cartan formula now proves the stability of every
$d_0^pD_k$.

Each nonzero term in \eqref{eq:milnor-action} raises the $d_0$-exponent
by exactly one.  Thus no monomial at exponent $p$ is hit.  Above this
floor, a monomial with some even $b_j$ is the $q_{j-1}$-image of the
monomial obtained by decreasing $c$ by one and increasing $b_j$ by one.
This preimage is still in $d_0^pD_k$.  If all $b_j$ are odd, the Euler
primitive on the monomial at exponent $c-1$ has coefficient
\[
 (c-1)+\sum_{j=1}^{k-1}b_j\equiv c+k\pmod2.
\]
It hits the given monomial exactly when $c\not\equiv k\pmod2$.
The images of the first $k$ primitives are spanned by these individual
monomials.  Later primitives add no new images: a square coefficient
passes inside any of the derivations.  This proves the basis assertion,
including the case $k=1$, when the list of $b_j$ is empty.
\end{proof}

\subsection{Cartier halving, with its exponent threshold}

We call the span of the monomials in \eqref{eq:coinvariants} the
\emph{above-floor part} of the coinvariants.  The following lemma keeps
track both of the action and of the first exponent in this part.

\begin{lemma}\label{lem:cartier}
The above-floor part of $H_0(E,d_0^pD_k)$ is an $A/K$-submodule.
After identifying $A/K$ with the Steenrod algebra by halving degrees,
the map
\begin{equation}\label{eq:cartier}
 C:\quad d_0^c\prod_{j=1}^{k-1}d_j^{b_j}
 \longmapsto
 d_0^{(c+k-2)/2}\prod_{j=1}^{k-1}d_j^{(b_j-1)/2}
\end{equation}
identifies it with an ideal in $D_k$.  More precisely, if
$c_0$ is the least integer greater than $p$ with $c_0\equiv k\pmod2$,
the target ideal is $d_0^{(c_0+k-2)/2}D_k$.
Internal polynomial degrees transform by $d\mapsto(d-k)/2$.
\end{lemma}

\begin{proof}
On $P_k$ define $C_x$ to be zero on a monomial unless all its exponents
are odd, and in the latter case set
\[
 C_x(x_1^{2a_1+1}\cdots x_k^{2a_k+1})
   =x_1^{a_1}\cdots x_k^{a_k}.
\]
This is the dual Kameko squaring map
\cite{Kameko1990}; see \cite[\S2, p.~4069]{Hung2005} for the
dual formula. It is also the polynomial top-form version of
Cartier's differential-form construction \cite{Cartier1957};
see also \cite[Chapter~II, \S6, pp.~200--203]{Cartier1958}.
We include the polynomial identities and naturality needed for its
restriction to the present ideals.
The binomial formula for squares on a one-variable polynomial, together
with the Cartan formula, gives
\begin{equation}\label{eq:algebra-cartier-action}
 C_x\Sq^{2i}=\Sq^iC_x,\qquad C_x\Sq^{2i+1}=0,
 \qquad C_x(a^2f)=aC_x(f).
\end{equation}
Indeed, on an odd exponent the parity of
$\binom{2a+1}{2i}$ is the parity of $\binom ai$; an odd square makes
that exponent even.  An even input exponent can never contribute an
odd output exponent, since the corresponding binomial coefficient then
vanishes.  Applying these facts in all variables proves the identities.

The operator $C_x$ commutes with $\GL_k(\F_2)$.  To see this, write a
polynomial uniquely as a sum of square-free monomials with square
coefficients.  The coefficient of $x_1\cdots x_k$ after a linear change
of variables is multiplied by its determinant; no square-free monomial
of smaller degree can contribute to that coefficient.  Over $\F_2$
the determinant is one, and taking the square root of the coefficient
commutes with the linear change of variables.  Hence $C_x$ preserves
$D_k$.

We compute this restriction by an elementary form of Cartier
naturality.  If $y_1,\ldots,y_k$ are algebraically independent
polynomials and $J=\det(\partial y_j/\partial x_i)$, then
\begin{equation}\label{eq:algebra-cartier-jacobian}
 C_x\bigl(Jf(y_1,\ldots,y_k)\bigr)
       =J\,C_y(f)(y_1,\ldots,y_k).
\end{equation}
Here $C_y$ is the same odd-exponent extractor in formal variables $y_j$.
For a proof, the polynomial differential-form complex in one variable
has cohomology generated by $1$ and $[x\,dx]$ over $\F_2[x^2]$.
Tensoring these complexes gives the corresponding assertion in $k$
variables.  Thus there is an isomorphism from the algebra of forms,
regarded without differential, to its cohomology, given by
\[
 f\longmapsto f^2,\qquad dx_i\longmapsto[x_i\,dx_i].
\]
This isomorphism is natural under polynomial substitution.  On
differentials the required fact is that $f\mapsto[f\,df]$ is additive
and satisfies the derivation rule with coefficients squared: its
additivity follows because the cross term is $d(fg)$, and its product
rule follows by expanding $(fg)d(fg)$.  In top degree the inverse
isomorphism sends $[h\,dx_1\cdots dx_k]$ to
$C_x(h)\,dx_1\cdots dx_k$.  Pulling back a top form in the $y$
variables proves \eqref{eq:algebra-cartier-jacobian}.

For Dickson coordinates $y_j=d_{j-1}$ the Jacobian is
\begin{equation}\label{eq:algebra-dickson-jacobian}
 J=d_0^{k-1}.
\end{equation}
In fact, differentiating $P(x_s)=0$ with respect to each $x_i$ gives
the matrix identity whose two factors have entries $x_s^{2^j}$ and
$\partial d_j/\partial x_i$, and whose product is $d_0$ times the
identity matrix.  The determinant of the first factor is $d_0$: it is
nonzero, is divisible by every nonzero binary linear form, and has
degree $2^k-1$.  Taking determinants gives
\eqref{eq:algebra-dickson-jacobian}.

For $c\geq k-1$, apply
\eqref{eq:algebra-cartier-jacobian} to
$f=d_0^{c-k+1}\prod_{j\geq1}d_j^{b_j}$.  It is zero unless
$c\equiv k\pmod2$ and every $b_j$ is odd; otherwise it is exactly
\eqref{eq:cartier}.  The same assertion holds for smaller nonnegative
$c$.  Multiply first by a sufficiently large even power $d_0^{2N}$,
use $C_x(d_0^{2N}f)=d_0^NC_x(f)$, and cancel $d_0^N$ in the polynomial
domain.  The output exponents in the nonzero cases are nonnegative.

The description of the hits in
Lemma~\ref{lem:milnor-coinvariants} now shows that $C_x$ descends to
coinvariants and is a bijection on the stated above-floor monomials.
Squares do not decrease the $d_0$-exponent: this follows from the
divisibility of the total square of $d_0$ used above and the Cartan
formula.  Consequently the above-floor span is a submodule.
Equation~\eqref{eq:algebra-cartier-action} gives the claimed action on
its image.  The least possible target exponent is the one displayed
in the statement, and a direct degree calculation in
\eqref{eq:cartier} gives the change $d\mapsto(d-k)/2$.
\end{proof}

\begin{remark}
On its image, the inverse monomial transformation in
\eqref{eq:cartier} is
\(f\mapsto f^2d_0^{2-k}d_1\cdots d_{k-1}\), with nonnegative
resulting exponents. Giambalvo and Peterson use this transformation
in their study of Dickson ideals \cite[\S4]{GiambalvoPeterson2001}.
Lemma~\ref{lem:cartier} specifies the simultaneous Milnor-coinvariant
quotient and the exact exponent floor used in the argument here.
\end{remark}

\subsection{The lifted functional and its exponent bound}

We first use Cartier halving on the functional from the preceding
level. Its ideal has exponent floor $n-1$, while the monomial it
detects already belongs to the next ideal, with floor $n$. The
resulting bound will control which ideal occurs when we halve the
tensor itself, including whether the floor can contribute.

\begin{lemma}\label{lem:algebra-functional-halving}
Let $n=e+1>0$, let $m$ be even, and let $v$ be a degree-$m$ monomial
in $M_k(n)$ with $d_0$-exponent $c$.  Suppose that its underlying
monomial is detected by an annihilated functional on $M_k(e)$ in
degree $m-1$.  If $b=2^{\nu_2(m)}$, then
\begin{equation}\label{eq:halving-bounds}
 n\equiv k\pmod b,\qquad c\geq n+b-2.
\end{equation}
If $m$ is a power of two, this functional hypothesis alone forces
$n=k=1$.
\end{lemma}

\begin{proof}
The monomial lies above the exponent floor $e$ in $M_k(e)$.
Since the functional factors through exterior coinvariants,
Lemma~\ref{lem:milnor-coinvariants} forces
$c\equiv k\pmod2$ and all its other exponents to be odd.
Its total degree there is $m-1$, so
\eqref{eq:algebra-degree-parity} gives $n\equiv k\pmod2$.

Restrict the induced functional to the above-floor submodule and use
Lemma~\ref{lem:cartier}.  The first allowed exponent in this submodule
is $c=n$, so its image has first exponent
\[
 e'=\frac{n+k-2}{2}.
\]
Give that image the grading shift $e'$, and put
\begin{equation}\label{eq:algebra-functional-recursion}
 m'=\frac m2,\qquad n'=\frac{n+k}{2},\qquad
 e'=n'-1,\qquad c'=\frac{c+k-2}{2}.
\end{equation}
The image is exactly $M_k(e')$, and its paired monomial has degree
$m'-1$: if the old internal degree is $d$, then
$e'+(d-k)/2=(m-2)/2$.
The new functional is again annihilated, because the Cartier map
intertwines the halved action.  If $m'$ is even, its paired monomial
cannot be at the floor $c'=e'$; a monomial at that floor has even
total degree, whereas $m'-1$ is odd.  It therefore lies in the next
ideal and satisfies the same hypothesis, allowing the step to repeat.

Iterate while the paired parameter $m$ is even, for
$t=\nu_2(m)$ steps.  With $b=2^t$, the resulting parameters obey
\[
 n_t=\frac{n+(b-1)k}{b},\qquad
 c_t-n_t+2=\frac{c-n+2}{b},\qquad c_t\geq n_t-1.
\]
The integrality of $n_t$ gives the congruence in
\eqref{eq:halving-bounds}, and the last two relations give its
inequality.  All shifts $e_t=n_t-1$ are nonnegative.

If $m=b$, the final paired monomial has degree zero in $M_k(e_t)$.
The smallest degree in that module is $2^ke_t$, so $e_t=0$.
Thus
\[
 n+(m-1)k=m.
\]
Since $n,k\geq1$, this equality forces $n=k=1$.
\end{proof}

\subsection{Injectivity for degree-one Tor}

The exact sequence $0\to I\to A\to\F_2\to0$ identifies
$\Tor_1^A(\F_2,M)$ with the kernel of the action map
$I\otimes_A M\to M$.  We refer to elements of this kernel as
\emph{cycles}.  A cycle represents zero precisely when it is zero in
the balanced tensor product.

\begin{lemma}\label{lem:tor-injection}
Let $n\geq0$ and $n\equiv k\pmod2$.  On cycles of even total degree,
the projection
\begin{equation}\label{eq:tor-injection}
 I\otimes_A M_k(n)\longrightarrow
 I_{A/K}\otimes_{A/K}H_0(E,M_k(n))
\end{equation}
is injective.
\end{lemma}

\begin{proof}
Put $M=M_k(n)$ and $B=A/K$.
The exact sequence of right $A$-modules
$K\to I\to I_B\to0$, followed by tensoring with $M$, shows that the
kernel of \eqref{eq:tor-injection} is the image of $K\otimes_A M$.
Here we have used
\[
 I_B\otimes_A M
 =I_B\otimes_B(B\otimes_A M)
 =I_B\otimes_BH_0(E,M).
\]
Multiplication gives a surjection $I_E\otimes_EA\to K=I_EA$;
tensoring again gives a surjection
$I_E\otimes_EM\to K\otimes_AM$.
Thus every element in this kernel has a preimage in
$I_E\otimes_EM$.  The action maps from both tensor products to $M$
commute with this map.  A preimage of a cycle is therefore a cycle
itself.

We prove that every even-total-degree cycle in $I_E\otimes_EM$ is
zero.  Write a homogeneous such cycle as
\[
 z=\sum_i q_i\otimes z_i.
\]
Each $z_i$ has odd total degree in $M$.  By
\eqref{eq:algebra-degree-parity}, its monomials have $d_0$-exponent
$c\geq n+1$ and parity opposite to $n$, hence opposite to $k$.
Lemma~\ref{lem:milnor-coinvariants} gives $z_i\in I_EM$.

This last property allows us to replace later primitives by the first
$k$ inside the tensor product.  If
$q_i=\sum_{j<k}a_j^2q_j$ on $M$, then
\begin{equation}\label{eq:algebra-hit-balancing}
 q_i\otimes z_i=\sum_{j<k}q_j\otimes a_j^2z_i
 \qquad(z_i\in I_EM).
\end{equation}
It suffices to check this for $z_i=q_sy$.  Balancing, commutativity of
the exterior generators, and their vanishing on squares give
\[
 \begin{aligned}
 q_i\otimes q_sy
 &=q_s\otimes q_iy
 =\sum_{j<k}q_s\otimes a_j^2q_jy\\
 &=\sum_{j<k}q_s\otimes q_j(a_j^2y)
 =\sum_{j<k}q_j\otimes a_j^2q_sy.
 \end{aligned}
\]
Thus $z$ has a cycle representative using only $q_j$, $j<k$, whose
coefficients still have exponents at least $n+1$.

Let $E_{[k]}=\Lambda(q_0,\ldots,q_{k-1})$ and
$L=D_k[d_0^{-1}]$.  We first show that $L$ is projective over
$E_{[k]}$.  For $T\subseteq\{0,\ldots,k-1\}$ write
$q_T=\prod_{i\in T}q_i$, and put
\[
 v_0=d_0^{-k}\prod_{j=1}^{k-1}d_j.
\]
The first $k-1$ primitives take $v_0$ to $d_0^{-1}$, and the last,
Euler primitive takes $d_0^{-1}$ to $1$.  Hence
$q_{\{0,\ldots,k-1\}}v_0=1$.
The action map from the free module $E_{[k]}\otimes_{\F_2}L$ to $L$
has the following section:
\begin{equation}\label{eq:algebra-projective-section}
 s(x)=\sum_{T\subseteq\{0,\ldots,k-1\}}
        q_T\otimes v_0\bigl(q_{T^c}x\bigr).
\end{equation}
It is $E_{[k]}$-linear: multiplying its first factor by $q_i$ and
reindexing the subsets gives $s(q_ix)$, since terms with repeated
$q_i$ vanish.  To compute the composite with the action map, expand
each $q_T(v_0q_{T^c}x)$ by the Cartan formula.  For a fixed subset $U$,
the term $(q_Uv_0)(q_{U^c}x)$ occurs once for every $T$ with
$U\subseteq T\subseteq\{0,\ldots,k-1\}$.  Its coefficient is
$2^{k-|U|}$, so only the full subset remains.  The composite is
therefore $x$.  The formula preserves degrees, since $v_0$ has the
negative of the degree of the product of all the $q_i$.
This proves the asserted projectivity.

Consider now the tensor product of the usual periodic free resolutions
for the one-generator exterior algebras.  After tensoring with an
$E_{[k]}$-module, its first differential is
\[
 (u_0,\ldots,u_{k-1})\longmapsto\sum_{i<k}q_iu_i.
\]
The degree-two terms have one generator for each repeated pair $i,i$
and each distinct pair $i,j$. Writing $e_i$ for the degree-one
generators, their images are
\[
 e_i\otimes q_i y,\qquad
 e_i\otimes q_jy+e_j\otimes q_iy\quad(i<j).
\]
The cokernel of this differential is
$I_{E_{[k]}}\otimes_{E_{[k]}}M$: the displayed relations become
$q_i^2\otimes y=0$ and
$(q_iq_j+q_jq_i)\otimes y=0$ under balancing.
Projectivity of $L$ implies that every degree-one
cycle in this complex with coefficients in $L$ has a degree-two
filling.

Apply this to the cycle representative obtained above, ignoring its
harmless shift by $n$.  Choose a filling with Laurent-polynomial
coefficients.  In each of the first two differentials, every nonzero
term raises the $d_0$-exponent of its coefficient by exactly one, by
\eqref{eq:milnor-action}.  Delete from the filling all monomials with
exponent less than $n$.  Their differential has exponents at most $n$,
whereas the retained part has exponents at least $n+1$.
Since the given cycle has only exponents at least $n+1$, the retained
part is still a filling.  Its coefficients belong to $d_0^nD_k$.
All Laurent polynomials and sums used here are finite, so this deletion
requires no completion.

We have proved that the representative is a boundary already over
$E_{[k]}$ with coefficients in $M$.  Equivalently, it is zero in
$I_{E_{[k]}}\otimes_{E_{[k]}}M$, and hence also in
$I_E\otimes_EM$.  Its image in $I\otimes_AM$ is zero, which proves
the lemma.
\end{proof}

\subsection{Proof of the obstruction lemma}

\begin{lemma}\label{lem:algebra-odd-tensor}
If $n>0$, $m\geq3$ is odd, and $v\in M_k(n)$ is homogeneous of
degree $m$, then $\Sq^m\otimes v=0$ in $I\otimes_AM_k(n)$.
\end{lemma}

\begin{proof}
The Adem relation $\Sq^m=\Sq^1\Sq^{m-1}$ gives
$\Sq^m\otimes v=\Sq^1\otimes\Sq^{m-1}v$.
For $n\geq2$ the second operation exceeds the internal polynomial
degree $m-n$, so this is zero.  For $n=1$ it is the internal square
$v^2$.  By linearity we may take
$v=d_0^c\prod_{j\geq1}d_j^{b_j}$, where $c\geq1$.
For $k\geq2$, formula~\eqref{eq:milnor-action} gives
\[
 v^2=\Sq^1\left(
 d_0^{2c-1}d_1^{2b_1+1}\prod_{j=2}^{k-1}d_j^{2b_j}
 \right).
\]
The preimage belongs to $d_0D_k$.  For $k=1$ use instead
$\Sq^1(d_0^{2c-1})=d_0^{2c}$.  Balancing once more and using
$(\Sq^1)^2=0$ proves the assertion.
\end{proof}

The even-degree argument below uses the following square formulas
only in its final case. These are the square-action formulas of Hai and Hung,
recorded in \cite[Theorem~B]{Hung1991}; the proof below uses the
root-polynomial method of \cite[\S II]{Wilkerson1983}.

\begin{lemma}\label{lem:algebra-near-top-square}
The total Steenrod square on a Dickson generator is
\begin{equation}\label{eq:algebra-total-dickson-square}
 \Sq(d_j)=d_j^2+
       \left(\sum_{i=0}^jd_i\right)
       \left(\sum_{a=j+1}^kd_a\right).
\end{equation}
In particular the operation one below the top on $d_j$ is
$d_0d_{j+1}$.  For a positive-degree monomial
$v=\prod_{j=0}^{k-1}d_j^{b_j}$ of internal degree $d$,
\begin{equation}\label{eq:algebra-one-below-top}
 \Sq^{d-1}v
 =v^2\sum_{j:\ b_j\ \text{odd}}\frac{d_0d_{j+1}}{d_j^2}.
\end{equation}
Every displayed summand is a polynomial.
\end{lemma}

\begin{proof}
Apply the total square simultaneously to the variables $x_s$ and $U$
in \eqref{eq:algebra-dickson-polynomial}.  Since
$\Sq(U+a)=(U+a)(U+a+1)$ and $P$ is additive,
\[
 \sum_{j=0}^k\Sq(d_j)(U+U^2)^{2^j}
 =P(U)P(U+1)=P(U)^2+P(1)P(U).
\]
Comparison of the coefficients of $U^{2^j}$ gives
$\Sq(d_0)=P(1)d_0$ and, for $1\leq j<k$,
\[
 \Sq(d_j)+\Sq(d_{j-1})=d_{j-1}^2+P(1)d_j.
\]
Summing these identities yields
\[
 \Sq(d_j)=P(1)\sum_{i=0}^jd_i+\sum_{i=0}^{j-1}d_i^2,
\]
which is \eqref{eq:algebra-total-dickson-square}.
The term of degree $2|d_j|-1$ in that formula is exactly
$d_0d_{j+1}$.  In the Cartan expansion of an operation one below the
top on a product, exactly one factor receives its operation one below
the top and every other factor receives its top operation.  Equal
choices cancel in pairs, leaving precisely
\eqref{eq:algebra-one-below-top}.
\end{proof}

\begin{lemma}\label{lem:algebra-even-vanishing}
Let $m$ be even and not a power of two, let $b=2^{\nu_2(m)}$, and
let $v\in M_k(n)$ be a degree-$m$ monomial with $n>0$ and
$d_0$-exponent $c$.  If \eqref{eq:halving-bounds} holds, then
$\Sq^m\otimes v=0$ in $I\otimes_AM_k(n)$.
\end{lemma}

\begin{proof}
The tensor is a cycle, because $\Sq^m$ exceeds the internal degree
$m-n$ of $v$.  We induct on $\nu_2(m)$, using
Lemma~\ref{lem:tor-injection} throughout; its parity assumption follows
from \eqref{eq:halving-bounds}. Here we halve coinvariants of the
current ideal, with floor $n$. This differs from the floor $n-1$
used for the functional in Lemma~\ref{lem:algebra-functional-halving}.

Suppose first that $b\geq4$.  The image of $v$ in exterior
coinvariants is either zero, in which case the injectivity lemma gives
the result immediately, or is an above-floor survivor.  Indeed
$c\geq n+b-2\geq n+2$.  In this case the first allowed exponent in
the above-floor part of $M_k(n)$ is $n+2$, and its Cartier image has
first exponent
\[
 n'=\frac{n+k}{2}.
\]
With this shift the image is $M_k(n')$, with total degrees halved.
The image of our cycle in \eqref{eq:tor-injection} consequently comes
from
$\Sq^{m/2}\otimes v'$ in $I\otimes_AM_k(n')$, where $v'=C(v)$.
Put $c'=(c+k-2)/2$ and $b'=b/2$.  Then
\[
 n'\equiv k\pmod{b'},\qquad
 c'-n'+2=\frac{c-n+2}{2}\geq b'.
\]
Thus the same bounds hold for $m/2,n',c'$.  Since $m/2$ is still
not a power of two, induction makes this latter tensor zero, and
injectivity gives the required vanishing of the original cycle.

It remains to treat $b=2$.  Write $m=2L$, with $L\geq3$ odd.
All coinvariants of $M_k(n)$ have even total degree: this is clear
at the floor from \eqref{eq:algebra-degree-parity}, and above it from
$c\equiv k\equiv n\pmod2$.  Regrade them by half and call the
resulting $A$-module $C_n$.  The projected tensor is
\[
 \Sq^L\otimes[v]
   =\Sq^1\otimes\Sq^{L-1}[v].
\]
The second factor is the class of the original operation
$\Sq^{m-2}v$.  It suffices to show that this class is zero or an
image of $\Sq^1$ in $C_n$, since balancing then introduces
$(\Sq^1)^2=0$.

If $n\geq3$, the operation already exceeds the internal degree of
$v$.  If $n=2$, then $k$ is even and at least two, and the operation
is the internal square $v^2$.  This lies above the floor and has all
nonzero-index exponents even, so it is zero in coinvariants by
Lemma~\ref{lem:milnor-coinvariants}.

Suppose $n=1$, so $k$ is odd.  For $k=1$ the coinvariant basis,
including the floor, consists of all odd monomials in $d_0$.
The map $C_x$ identifies $C_1$ with $\Sigma D_1$, and the image of
$v=d_0^{2L-1}$ is $d_0^{L-1}$.  Its image under $\Sq^{L-1}$ is
\[
 d_0^{2L-2}=\Sq^1(d_0^{2L-3}),
\]
as required.

Finally, suppose $n=1$ and $k\geq3$ is odd.  The internal degree of
$v$ is $m-1$, so Lemma~\ref{lem:algebra-near-top-square} applies to
$\Sq^{m-2}v$.  In each summand of
\eqref{eq:algebra-one-below-top}, at most one exponent of
$d_j$ with $j\geq1$ is odd.  Since there are at least two such
indices, every term above the floor dies in exterior coinvariants.
The only possible floor term comes from the index $j=0$ when
$v=d_0w$, where $w$ involves only $d_j$, $j\geq1$.  It is
$d_0d_1w^2$.

This remaining term is an image for the halved first square.  Formula
\eqref{eq:algebra-total-dickson-square} gives
$\Sq^2d_2=d_1$ and $\Sq^1d_0=\Sq^2d_0=0$ for $k\geq3$.
Together with \eqref{eq:milnor-action} and the Cartan formula, it
therefore gives the explicit identity
\begin{equation}\label{eq:algebra-floor-witness}
 \Sq^2(d_0d_2w^2)
 =d_0d_1w^2+
   d_0^3d_2\left(\frac{\partial w}{\partial d_1}\right)^2.
\end{equation}
The extra term is an exterior hit in $d_0D_k$, with preimage
\[
 q_0\left(
 d_0^2d_1d_2\left(\frac{\partial w}{\partial d_1}\right)^2
 \right)
 =d_0^3d_2\left(\frac{\partial w}{\partial d_1}\right)^2.
\]
Hence the class of $d_0d_1w^2$ is the original $\Sq^2$-image of
$[d_0d_2w^2]$, or equivalently a $\Sq^1$-image after halving.
This completes the final case, and the injectivity lemma finishes
the proof.
\end{proof}

\begin{proof}[Proof of Lemma~\ref{lem:algebra-obstruction}]
The odd assertion is Lemma~\ref{lem:algebra-odd-tensor}.
For even $m$, the functional hypothesis gives
\eqref{eq:halving-bounds} by
Lemma~\ref{lem:algebra-functional-halving}.  If $m$ is not a power
of two, Lemma~\ref{lem:algebra-even-vanishing} contradicts
\eqref{eq:obstruction-tensor}.  If $m$ is a power of two, the last
assertion of Lemma~\ref{lem:algebra-functional-halving} gives
$n=k=1$.
\end{proof}
